\section{The signed Tait partition function in the repository convention}
\label{sec:potts}

\subsection{Input, embedding, and signs}

Let $D$ be a classical one-component knot diagram with $n>0$ crossings.
Its shadow is a connected plane four-valent multigraph. The public input
validator checks that the supplied planar-diagram code has the required
edge incidences, one link component, and spherical rotation system.
The following construction uses that embedding, not an abstract graph
with its rotation data discarded.

Choose one checkerboard color. The Tait graph $G=(V,E)$ has one vertex for
each region of that color and one edge for each crossing. It can contain
loops and parallel edges. Write $v=|V|$. At a crossing $e$, designate as the
\emph{omitted} smoothing the one that leaves the two selected-color
corners disconnected across that crossing. Its bracket coefficient is
$A^{\epsilon_e}$, where $\epsilon_e\in\{1,-1\}$.
The included smoothing has coefficient $A^{-\epsilon_e}$.
This convention, which is checked from the local darts in the implementation,
avoids assuming that an alternating-diagram sign rule applies to all
diagrams. Put
\[
 s=\sum_{e\in E}\epsilon_e,\qquad \delta=-A^2-A^{-2}.
\]
We use the \emph{full} bracket $B_D(A)$, for which an isolated circle has
value $\delta$. The normalized Jones polynomial is
\begin{equation}
 V_D(A^{-4})=\frac{B_D(A)}{\delta(-A^3)^{\wr(D)}}.
 \label{eq:jones-normalization}
\end{equation}
Consequently the full bracket of an unknot diagram with writhe $\wr(D)$ is
$\delta(-A^3)^{\wr(D)}$. Losing the initial circle factor would invalidate the
obstruction.

\begin{lemma}[State circles]\label{lem:circles}
For $S\subseteq E$, let $k(S)$ be the number of components of the spanning
subgraph $(V,S)$, including isolated vertices. The associated smoothing state
has $c(S)=2k(S)-v+|S|$ circles.
\end{lemma}
\begin{proof}
A regular neighborhood of $(V,S)$ in the sphere retracts to that subgraph,
so its Euler characteristic is $v-|S|$. Its $k(S)$ components are planar
surfaces. If their total number of boundary components is $c$, their Euler
characteristic is $2k(S)-c$. Those boundary components are the smoothing
circles. Equating the two expressions proves the formula. An isolated vertex
has a disk neighborhood and contributes one circle, as required. Loops and
parallel edges are included in the same argument.
\end{proof}

\subsection{From smoothing states to spins}

The multivariate random-cluster partition function is
\[
 Z_G(Q,\{z_e\})=\sum_{S\subseteq E}Q^{k(S)}\prod_{e\in S}z_e.
\]
At a positive integer $Q=q$, the Fortuin--Kasteleyn identity is
\begin{equation}
 Z_G(q,\{z_e\})
 =\sum_{\sigma:V\to[q]}
   \prod_{e=uv}\bigl(1+z_e[\sigma(u)=\sigma(v)]\bigr).
 \label{eq:fk}
\end{equation}
To verify it, expand the product on the right. For each selected subset
$S$, the spin assignment must be constant on each component of $(V,S)$,
giving exactly $q^{k(S)}$ choices. This also proves the identity for loops,
parallel edges and isolated vertices \cite{sokal2005}.

\begin{proposition}[Signed partition-function formula]\label{prop:potts}
In the rational-function field, or after localizing at $\delta$, with
$Q=\delta^2$ and $z_e=\delta A^{-2\epsilon_e}$,
\begin{equation}
 B_D(A)=\delta^{-v}A^s Z_G(Q,\{z_e\}).
 \label{eq:full-bracket}
\end{equation}
\end{proposition}
\begin{proof}
The smoothing indexed by $S$ contributes
\[
 A^s\prod_{e\in S}A^{-2\epsilon_e}\,
 \delta^{2k(S)-v+|S|}.
\]
Factor out $\delta^{-v}A^s$ and sum. The remaining term for $S$ is
$(\delta^2)^{k(S)}\prod_{e\in S}(\delta A^{-2\epsilon_e})$.
\end{proof}

Set $x=A^4$. Then
\[
 Q=x+2+x^{-1},\qquad z_e=-1-x^{-\epsilon_e}.
\]
Thus the spin interaction has the especially simple form
\begin{equation}
 h_e(a,b)=
 \begin{cases}
  -x^{-\epsilon_e},&a=b,\\
  1,&a\ne b.
 \end{cases}
 \label{eq:interaction}
\end{equation}
This is the signed interaction used in the earlier tensor-network
formulation \cite{meichanetzidis2019}. It allows all crossings to be processed
as local equality tests, without tracking smoothing connectivity.

\subsection{Eliminating fourth roots from normalization}

Since $\delta=-(x+1)A^{-2}$, substituting
\eqref{eq:full-bracket} into \eqref{eq:jones-normalization} gives
\begin{equation}
 V_D(A^{-4})
 =(-1)^{\wr(D)+v+1}A^E
   \frac{Z_G}{(x+1)^{v+1}},
 \qquad E=s-3\wr(D)+2(v+1).
 \label{eq:normalize-E}
\end{equation}

\begin{lemma}[Integral normalization exponent]\label{lem:parity}
For a knot, $E$ is divisible by four.
\end{lemma}
\begin{proof}
For a smoothing state with coefficient exponent $s_\sigma$ and circle
count $c_\sigma$, the residue
\[
 s_\sigma-3\wr(D)+2(c_\sigma-1)\pmod4
\]
is unchanged when one smoothing is flipped: the exponent changes by
$\pm2$ and the circle count by $\pm1$. In the oriented Seifert state,
$s_\sigma=\wr(D)$. The genus formula for a one-component Seifert surface
gives $c_\sigma-1\equiv n\equiv\wr(D)\pmod2$, so the residue is zero.
In the all-omitted state, $s_\sigma=s$ and $c_\sigma=v$.
Its expression differs from $E$ by four.
\end{proof}

Writing $h=E/4$, the exact equality test can therefore be performed as
\begin{equation}
 Z_G\stackrel{?}{=}Z_U,\qquad
 Z_U=(-1)^{\wr(D)+v+1}x^{-h}(x+1)^{v+1}.
 \label{eq:unknot-partition}
\end{equation}
No fourth root of $x$ is needed. This is an equality of algebraic integers in
the ring introduced in Section~\ref{sec:exact}. The comparison is obtained by
multiplying through the nonzero normalization factor; the implementation
never needs to divide by $x+1$.

The empty crossing diagram is handled separately as a single circle, with
$Z_G=Z_U=q$. For a link with $\ell$ components the parity residue in the lemma
would instead be $2(\ell-1)\pmod4$. This article and the public recognizer
require a knot; accepting all links without revisiting normalization would
be an unsupported generalization.
